\documentclass[12pt]{article}
\usepackage[margin=0.75in]{geometry}
\usepackage{fontspec}
\setmainfont{Latin Modern Roman}
\usepackage{microtype}
\usepackage{fancyhdr}
\usepackage{titlesec}
\usepackage{enumitem}
\usepackage{xcolor}
\usepackage[hidelinks]{hyperref}
\setlength{\parindent}{1.1em}
\setlength{\parskip}{0.38em}
\setlength{\headheight}{15pt}
\titleformat{\section}{\large\bfseries\scshape}{}{0pt}{}
\titleformat{\subsection}{\normalsize\bfseries}{}{0pt}{}
\pagestyle{fancy}
\fancyhf{}
\fancyfoot[C]{\thepage}
\renewcommand{\headrulewidth}{0.4pt}
\newcommand{\focusbox}[1]{\par\vspace{0.4em}\noindent\begingroup\setlength{\fboxsep}{6pt}\colorbox{gray!12}{\begin{minipage}{\dimexpr\linewidth-2\fboxsep\relax}#1\end{minipage}}\endgroup\par\vspace{0.5em}}
\newcommand{\studentline}{\vspace{0.35em}\hrule\vspace{0.65em}}
\newcommand{\shortline}{\vspace{0.25em}\hrule\vspace{0.45em}}

\fancyhead[L]{Whately: Validity and Truth Lab}
\fancyhead[R]{Student Handout}
\begin{document}
\begin{center}
{\LARGE\bfseries Week 6 Argument Lab}\par\vspace{0.25em}
{\large\scshape Good Form, Bad Facts}\par\vspace{0.6em}
\rule{0.82\textwidth}{0.4pt}
\end{center}

\section*{Name: \rule{2.4in}{0.4pt} \hfill Date: \rule{1.3in}{0.4pt}}

\section*{Purpose}
This lab turns Week 6 reading into a portfolio artifact. You will diagnose arguments with two columns: a \textbf{form test} and a \textbf{fact test}. The goal is to stop using the word \textit{valid} when you really mean \textit{true}, \textit{persuasive}, \textit{likely}, or \textit{morally attractive}.

\section*{Part I: Reference Box}
\focusbox{\textbf{Form test:} Does the conclusion follow from the premises? \textbf{Fact test:} Are the premises true or adequately supported? \textbf{Sound argument:} valid form plus true premises. \textbf{Unsound argument:} invalid form or at least one false premise. \textbf{Not yet proved:} form may be valid, but a key premise remains doubtful or unsupported.}

\section*{Part II: Model Diagnosis}
\textbf{Argument:}
\begin{itemize}[leftmargin=*]
\item All citizens who knowingly sell public secrets for private money betray their city.
\item Marcus knowingly sold public secrets for private money.
\item Therefore, Marcus betrayed his city.
\end{itemize}

\begin{itemize}[leftmargin=*]
\item \textbf{Form test:} Valid. If both premises are true, the conclusion follows.
\item \textbf{Fact test:} The second premise is the pressure point. Did Marcus knowingly sell secrets, and was it for private money?
\item \textbf{Diagnosis:} The argument is sound only if that factual premise is proved. Until then, the conclusion is not yet proved.
\end{itemize}

\section*{Part III: Diagnose the Arguments}
For each argument, complete both tests. Be precise: do not call an argument invalid merely because you doubt a premise.

\subsection*{1. Valid Form, False Premise}
\textbf{Argument:} All birds can breathe underwater. All eagles are birds. Therefore, all eagles can breathe underwater.

\textbf{Form test: Does the conclusion follow from the premises?}\studentline
\textbf{Fact test: Which premise is false or doubtful?}\studentline
\textbf{Final diagnosis: sound, unsound, or not yet proved?}\studentline

\subsection*{2. True Conclusion, Bad Proof}
\textbf{Argument:} All careful students own notebooks. Elena owns a notebook. Therefore, Elena is a careful student.

\textbf{Form test: Does the conclusion follow from the premises?}\studentline
\textbf{Could the conclusion still be true? Why does that not make this proof good?}\studentline\studentline
\textbf{Final diagnosis:}\studentline

\subsection*{3. Valid Form, Doubtful Premise}
\textbf{Argument:} All speeches that intentionally deceive the people are unjust. This speech intentionally deceives the people. Therefore, this speech is unjust.

\textbf{Form test:}\studentline
\textbf{Fact test: What would need to be proved about the speech?}\studentline\studentline
\textbf{Final diagnosis:}\studentline

\section*{Part IV: Two-Column Argument Diagnosis}
Choose one argument from Part III and fill out the diagnosis chart.

\begin{itemize}[leftmargin=*]
\item \textbf{Conclusion:}\studentline
\item \textbf{Form judgment: valid or invalid?}\studentline
\item \textbf{Premise pressure point:}\studentline
\item \textbf{Fact judgment: true, false, doubtful, or not yet established?}\studentline
\item \textbf{What evidence would settle the fact test?}\studentline\studentline
\item \textbf{Final diagnosis:}\studentline
\end{itemize}

\section*{Part V: Julius Caesar Application}
Brutus's public argument can be reconstructed this way:
\begin{itemize}[leftmargin=*]
\item All men whose ambition would enslave Rome deserve death for Rome's freedom.
\item Caesar was a man whose ambition would enslave Rome.
\item Therefore, Caesar deserved death for Rome's freedom.
\end{itemize}

\textbf{Form test: Is this reconstruction valid? Explain.}\studentline\studentline
\textbf{Fact test: Which premise does Antony attack?}\studentline
\textbf{What evidence does Antony offer against that premise?}\studentline\studentline
\textbf{Does Antony prove Brutus wrong, or does he make Brutus's argument not yet proved? Explain carefully.}\studentline\studentline

\section*{Part VI: Repair or Narrow}
Choose one weak argument from this lab. Either strengthen a doubtful premise with evidence or narrow the conclusion so it no longer overclaims.

\textbf{Original argument:}\studentline
\textbf{Problem: invalid form, false premise, doubtful premise, or overclaim?}\studentline
\textbf{Repaired or narrowed argument:}\studentline\studentline\studentline

\section*{Part VII: Portfolio Artifact}
Create a one-page \textbf{Form Test / Fact Test Diagnosis} for your Logic Portfolio. It should include:
\begin{enumerate}[leftmargin=*]
\item a real or literary argument;
\item its conclusion;
\item a form judgment;
\item a fact judgment;
\item the premise pressure point;
\item a final diagnosis: sound, unsound, or not yet proved;
\item one sentence explaining why validity alone is not enough.
\end{enumerate}

\focusbox{\textbf{Portfolio summary starter:} The argument's form is \rule{1.4in}{0.4pt}, but the premise that most needs testing is \rule{1.8in}{0.4pt}. This shows that validity is not enough because \rule{2.0in}{0.4pt}.}

\vspace{1.0in}
\end{document}
